\documentclass[a4paper]{esannV2}
\input{common/macros.tex}
% Disable margin protrusion to keep punctuation inside the ESANN printing area.
\usepackage[protrusion=false]{microtype}
\usepackage[hidelinks]{hyperref}
% Stay within the centered official printing area, allowing for font overhangs.
\setlength{\paperwidth}{21cm}\setlength{\paperheight}{29.7cm}
\setlength{\textwidth}{12.1cm}\setlength{\textheight}{19.0cm}
\setlength{\oddsidemargin}{1.91cm}\setlength{\evensidemargin}{1.91cm}
\setlength{\topmargin}{2.76cm}\setlength{\headheight}{0cm}\setlength{\headsep}{0cm}
\ifdefined\pdfpagewidth
  \pdfpagewidth=\paperwidth\pdfpageheight=\paperheight
\else
  \special{papersize=210mm,297mm}
\fi
\ifdefined\pdfimageresolution\pdfimageresolution=600\fi
\interfootnotelinepenalty=10000
\brokenpenalty=10000
\newtheorem{proposition}{Proposition}
\hypersetup{pdftitle={Building Activity Representations with Pascal-Coded Signatures},pdfauthor={Erlend Aune; Fabian N. Harang},pdfsubject={ESANN 2027}}

\begin{document}
\title{Building Activity Representations with\\Pascal-Coded Signatures}
\author{Erlend Aune$^{1}$ \quad Fabian N. Harang$^{2}$\\
$^{1}$Cognite; Norwegian University of Science and Technology (NTNU), Trondheim\\
$^{2}$BI Norwegian Business School;\\
Centre for Sustainable, Risk-Averse and Ethical Artificial Intelligence\\
(SURE-AI), Simula Research Laboratory}
\maketitle
\begin{abstract}
Movement timing can distinguish activities with the same displacement.
We study a polynomial
lift of path increments that exposes selected positional interactions as
endpoints and signed areas. Expressed in a Pascal basis, these coordinates admit
explicit temporal and spectral interpretations. We explain their relation to
ordinary signatures, distinguish signature depth from feature dimension, and construct
features that retain sensor levels, movement and rhythm. In exploratory
experiments, 368 features with a linear ridge classifier achieve 98.03\% accuracy
on UCI HAR's official split: smartphone motion recordings of six everyday
activities, with different participants for training and testing.
Conventional-signature controls and a
landmark gesture example clarify the roles of compactness, channel interactions
and representation choice. Code and extended technical notes can be found at:\\
\url{https://github.com/Froskekongen/pascal_coded_application}.
\end{abstract}

\section{From activity observations to path features}
Activity recognition can depend on movement timing even when overall displacement
is unchanged. The two hand-height paths
in Figure~\ref{fig:movement} start at zero, finish at two, and have the same
average height, but place their rises at the interval ends or middle. A useful
representation should expose this distinction to a simple classifier.

\begin{figure}[!ht]
\centering
\includegraphics[width=.98\linewidth]{evidence/natops/movement_schedules.pdf}
\caption{Two unit rises with the same initial, final and average heights.
Their spacing distinguishes the two schedules.}
\label{fig:movement}
\end{figure}

Path signatures describe sequential data through iterated integrals.
Their use in activity recognition depends on how observations are turned into
paths and which interactions are retained~\cite{yang,generalised}.
FRUITS weights iterated sums
by index separation~\cite{fruits}, while Volterra signatures admit exponential,
fractional and polynomial memory kernels~\cite{volterra}.
We study a fixed polynomial lift of discrete increments, expressed in a
Pascal basis. Our contribution is the interpretation and use of its depth-two
coordinates, together with observation designs that expose posture, movement
and rhythm to a linear readout.

\section{Polynomial coding at depth two}
Let $z_0,\ldots,z_{R-1}\in\R^C$, with $R\geq2$,
be the increments of a piecewise-linear path; for observed positions, $z_r=x_{r+1}-x_r$.
The first two signature levels are
\begin{equation}
 s_i=\sum_r z_r^i,\qquad
 B_{ij}=\sum_{r<t}z_r^i z_t^j+\frac12\sum_r z_r^i z_r^j.
 \label{eq:signature}
\end{equation}
The endpoint $s$ records net displacement, while $B$ records ordered
interactions. Since $B_{ij}+B_{ji}=s_i s_j$, the order-sensitive part is the
signed area
\begin{equation}
 J_{ij}=\frac12(B_{ij}-B_{ji})
       =\frac12\sum_{r<t}(z_r^i z_t^j-z_r^j z_t^i),\qquad
 B_{ij}=\frac12s_i s_j+J_{ij}.
 \label{eq:area}
\end{equation}
The area is closed by the endpoint chord. Endpoints and areas for $i<j$
form the depth-two logsignature.
Although $(s,B)$ and $(s,J)$ determine each other, a linear readout of the
former also has access to quadratic endpoint products.

These sums record order but do not explicitly weight increment separation.
An unclocked spatial signature cannot distinguish timings of the same geometric
traversal~\cite{generalised}. A time channel captures timing in the full signature,
but low-depth truncations can still merge differently timed movements. We
therefore construct a path whose endpoints summarize increment positions and
whose selected areas weight their separation.

\paragraph{Making position visible.}
We attach $P$ polynomially weighted copies to each increment,
for $1\leq P\leq R$:
\begin{equation}
 b_p(r)=\binom rp\binom{R-1}p^{-1},\qquad
 v_r=(b_p(r)z_r)_{p=0}^{P-1},\qquad 0\leq r<R,
 \label{eq:lift}
\end{equation}
where $\binom rp=0$ for $r<p$. The signature of the path accumulated
from $v_r$ is the \emph{Pascal-coded signature}: $b_0(r)=1$ retains the
original increments, $b_1(r)=r/(R-1)$ adds a linear position weight with values
between 0 and 1, and further copies add higher degrees. The weights act on
increments, not positions.

\paragraph{A distinction between two schedules.}
The one-dimensional paths in Figure~\ref{fig:movement} have increments
$z=(z_r)_{r=0}^{3}$: $(1,0,0,1)$ (spread) or $(0,1,1,0)$ (central).
Clocked increments $\tilde z_r=(1/4,z_r)$ give both paths endpoint $(1,2)$
and zero signed area by~\eqref{eq:area}.

For $R=4$, the weights \mbox{$u_r:=r/(R-1)$} satisfy
$(u_0,u_1,u_2,u_3)=(0,\tfrac13,\tfrac23,1)$.
Accumulating $v_r=(z_r,u_rz_r)\in\R^2$ gives both paths endpoint $(2,1)$
(total and position-weighted rise). By~\eqref{eq:area}, $J_{01}$ is the signed
area between coordinates $0$ (original) and $1$ (weighted):
\begin{equation}
 J_{01}=\frac12\sum_{r<t}(u_t-u_r)z_rz_t.
 \label{eq:sep}
\end{equation}
The rise separations $1-0=1$ and $\tfrac23-\tfrac13=\tfrac13$ give areas
$1/2$ and $1/6$. A linear readout separates the schedules at depth two;
the clocked paths first differ at depth three.

Both representations use two channels. Weighted second-order interactions
expose spacing and can be easier to interpret than higher-order terms at
comparable cost.
No general depth bound follows; in full coding, $P$ copies increase $C$
channels to $PC$, so lower depth alone implies no computational saving.

The binomial basis makes index shifts lower-triangular via Vandermonde's
identity, so locally coded segment signatures concatenate without revisiting
raw steps. At fixed $R$, it spans the same space as monomials of degree below
$P$. The resulting signature is also a Volterra signature~\cite{volterra} with
a position-only kernel held constant on each increment interval.

\section{What should the increments represent?}
The increments determine what information the features retain. For body
landmarks, differences between successive positions record movement but lose the
starting pose. We include this pose by adding a step from a fixed origin to the
first position, as in our gesture experiment. For inertial sensors, steady
readings can indicate posture. Taking differences would remove these steady
levels. We instead use scaled readings as increments, $z_r:=a_r/R$, so the
endpoint $s:=R^{-1}\sum_r a_r$ is the average reading. Additional polynomially
weighted copies add timing information while retaining this average.

For rhythm, Fourier analysis separates oscillations by frequency (how quickly
they repeat). With our normalization, total Fourier power
is the average squared reading, a measure of signal strength.
For nonzero signals, dividing each frequency's power by the total gives its energy
share $p_k\geq0$, with $\sum_k p_k=1$. We order $K\geq2$ equally spaced
nonnegative frequencies from low to high, using the normalized grid
$u_k:=k/(K-1)$. Zero frequency represents a constant level. We apply the Pascal
lift~\eqref{eq:lift} to the shares as increments, $z_k:=p_k$.

The first endpoint coordinate is their sum, $1$. The second is the
so-called spectral centroid $s_1:=\sum_k u_kp_k$: the energy-weighted average frequency
on this grid. The signed area describes how widely the energy is spread:
\begin{equation}
 J_{01}=\frac12\sum_{k<l}(u_l-u_k)p_kp_l
       =\frac14\mathbb E|U-V|,
 \label{eq:spectrum}
\end{equation}
where $U,V$ are independent frequency draws with $\Pr(U=u_k)=p_k$.
Their average absolute separation, $\mathbb E|U-V|$, is the Gini mean
difference~\cite{gini}.

Splitting a peak into two equal peaks, equally far above
and below it, preserves the centroid but increases $J_{01}$. Power ignores
where oscillations start in their cycle (phase). Normalization also removes
overall strength, which we retain through separate amplitude summaries.

An ordinary signature recovers the same summaries from a two-channel staircase
$(F,M)$: $F$ advances through $u_k$; $M$ accumulates power shares while
$F$ is fixed. Here $S^{MFM}=2J_{01}$. At fixed $K\geq3$, the scalar $P=3$
Pascal endpoints and areas are fixed linear combinations of staircase tensor
coordinates through depth five. The companion gives the derivation from
polynomial one-form signature calculus~\cite[Lem.~4.1]{polynomialoneforms}.

\section{A representation for phone activities}
We build \emph{compact368} from three feature blocks to recognize
UCI HAR's six activities from smartphone sensors (50\,Hz, 128-sample
windows)~\cite{har}. To improve conditioning, we weight increments by $Q_{rp}$
from $W=QH$, $W_{rp}=b_p(r)$, with $Q$ orthonormal and $H$ upper triangular
with positive diagonal. This QR gives identical $Q$ for ascending Pascal or
monomial columns in exact arithmetic.

\paragraph{Sensor levels and temporal interactions: 171 features.}
To retain posture, we standardize total acceleration and gyroscope (three axes
each) using training-sample means and population standard deviations, without
centering windows separately. Readings divided by 128 become increments. Three copies
give 18 channels and $18+\binom{18}{2}=171$ endpoints and areas. The
constant-copy endpoint is the standardized window mean divided by
$\sqrt{128}$, so static posture remains represented.

\paragraph{Spectral shape and scale: 140 features.}
We use five triads: body acceleration,
gravity (total minus body acceleration),
gyroscope, and first differences of body acceleration and gyroscope.
After training-set axis standardization, we append each triad's norm to obtain
20 scalar signals. Each normalized one-sided power spectrum contributes five
nonconstant coordinates from its three-copy logsignature; the constant
mass endpoint is removed. Means and population standard deviations of
unnormalized bin amplitudes add 40 scale features.

\paragraph{Strength, persistence and alignment: 57 features.}
We scale body acceleration, gravity and gyroscope by their training RMS.
Twelve streams (three squared norms, three simultaneous pairwise dot products
and six self-dot products at 40 and 100\,ms lags) contribute means, standard
deviations and skewness. Three pairwise cosine-similarity streams add means
and standard deviations. Standardized linear trends are retained for all fifteen streams.

After training-only feature standardization and constant-column removal, these
features feed a one-versus-rest ridge classifier with targets $+1/-1$, an
unpenalized intercept and largest-score prediction.

\section{Empirical illustrations and their scope}

We replace each coded branch of \emph{compact368} in turn with a conventional
signature (Table~\ref{tab:har}). Both studies are exploratory because these
feature families grew out of iterative work on public benchmarks. On the official UCI
HAR split (7,352 training windows from 21 subjects; 2,947 test windows from
nine other subjects), the replacement controls come from a 187-configuration
comparison. Every row uses training-only input and feature scaling with a fixed
ridge penalty $\alpha=10$.


\begin{table}[!ht]
\centering
\begin{tabular}{@{}lrrr@{}}
\toprule
Representation & Features & Accuracy & Errors\\
\midrule
\emph{compact368} &368&98.03\%&58\\
Raw branch: lead--lag tensor, depth 2 &377&97.46\%&75\\
Spectral branch: staircase tensor, depth 4 &708&98.00\%&59\\
Supplied UCI HAR features &561&95.72\%&126\\
\bottomrule
\end{tabular}
\caption{Exploratory HAR comparisons on the official subject-disjoint split.
Each replacement retains the other \emph{compact368} blocks; counts are retained
columns.}
\label{tab:har}
\end{table}

The raw control uses a clock-augmented, origin-anchored lead--lag path;
the staircase spectral control has more coordinates but only one extra error.

\paragraph{Landmark trajectories.}

On NATOPS~\cite{natopsoriginal,natops}, we compare channel interactions using depth-two
logsignatures of six gestures
tracked across eight 3D landmarks over 51 frames (180 training and 180 test
sequences). Appending a clock and prepending an origin anchor yields a
25-channel path, which $P=2$ coding~\eqref{eq:lift} lifts to 50 channels (Table~\ref{tab:natops}).
The \emph{full} regime retains all pairwise areas, whereas the \emph{local}
regime restricts areas to clock--channel and within-landmark coordinate pairs
across copies, plus the 25 same-channel copy areas; intermediate controls isolate
weighted endpoints and coded-clock areas. For each map, $\alpha\in\{10^k\}_{k=-6}^6$
is selected via shared five-fold stratified CV with in-fold standardization before
refitting on all 180 training sequences.

\begin{table}[!ht]
\centering
\begin{tabular}{@{}lrrrr@{}}
\toprule
Representation&Features&CV (\%)&Correct&Test (\%)\\
\midrule
Ordinary full&325&92.2&163/180&90.6\\
Moments + ordinary full$^*$&350&93.3&166/180&92.2\\
Coded-clock control (full)$^*$&423&92.2&163/180&90.6\\
Pascal full&1,275&92.2&173/180&96.1\\
\midrule
Ordinary local&73&92.8&167/180&92.8\\
Moments + ordinary local&98&92.8&165/180&91.7\\
Coded-clock control (local)&171&93.9&168/180&93.3\\
Pascal local&267&93.9&165/180&91.7\\
\bottomrule
\end{tabular}
\caption{NATOPS with shared training CV. Nominal counts include constant and
dependent coordinates. Moments add weighted endpoints; coded-clock controls
also add areas involving a clock copy. $^*$Controls added after the initial
seven outcomes.}
\label{tab:natops}
\end{table}

Coding improves the full map on test, though not the local map or CV (both
full maps select $\alpha=10$ at 92.2\% CV, below the local maps), correcting
12 ordinary-full errors while introducing two. This retrospective contrast
suggests useful coded spatial interactions, without isolating cross-landmark
effects from width and regularization. Computations use
\texttt{tensorsig}~\cite{tensorsig}.

\section{Conclusions and future work}
Pascal-coded logsignatures capture posture, timing and rhythm for linear readouts, with encouraging
exploratory accuracy on HAR and gesture trajectories. Future work includes
pre-registered subject-disjoint benchmarks, structured channel-pair
selection and learned positional bases to limit dimension.



\input{common/ai_disclosure.tex}
% The official ESANN example sets the bibliography in footnotesize.
\begin{footnotesize}
\begin{thebibliography}{99}
\bibitem{yang} W. Yang, T. Lyons, H. Ni, C. Schmid, and L. Jin. Developing the path signature methodology and its application to landmark-based human action recognition. In \href{https://doi.org/10.1007/978-3-030-98519-6_18}{\emph{Stochastic Analysis, Filtering, and Stochastic Optimization}}, pp.~431--464, Springer, 2022.
\bibitem{generalised} J. Morrill, A. Fermanian, P. Kidger, and T. Lyons. A generalised signature method for multivariate time series feature extraction. \href{https://arxiv.org/abs/2006.00873v2}{\emph{arXiv:2006.00873v2}}, 2021.
\bibitem{fruits} J. Diehl and R. Krieg. FRUITS: feature extraction using iterated sums for time series classification. \href{https://doi.org/10.1007/s10618-024-01068-1}{\emph{Data Mining and Knowledge Discovery}}, 38:4122--4156, 2024.
\bibitem{volterra} P. P. Hager, F. N. Harang, L. Pelizzari, and S. Tindel. The Volterra signature. \href{https://arxiv.org/abs/2603.04525v1}{\emph{arXiv:2603.04525v1}}, 2026.
\bibitem{gini} G. A. Koshevoy and K. Mosler. Multivariate Gini indices. \href{https://wisostat.uni-koeln.de/sites/dep_econometrics/pdf_publikationen/1997_KoshevoyM97a_Gini.pdf}{\emph{Journal of Multivariate Analysis}}, 60:252--276, 1997.
\bibitem{polynomialoneforms} H. Boedihardjo and X. Geng. The uniqueness of signature problem in the non-Markov setting. \href{https://arxiv.org/abs/1401.6165v2}{\emph{arXiv:1401.6165v2}}, 2014.
\bibitem{har} D. Anguita, A. Ghio, L. Oneto, X. Parra, and J. L. Reyes-Ortiz. A public domain dataset for human activity recognition using smartphones. \emph{ESANN}, pp.~437--442, 2013.
\bibitem{natopsoriginal} N. Ghouaiel, P.-F. Marteau, and M. Dupont. Continuous pattern detection and recognition in stream---a benchmark for online gesture recognition. \href{https://doi.org/10.1504/IJAPR.2017.085315}{\emph{International Journal of Applied Pattern Recognition}}, 4(2):146--160, 2017.
\bibitem{natops} TSML Archive Maintainers. NATOPS dataset. UEA time series classification archive. \href{https://doi.org/10.5281/zenodo.11206248}{Zenodo record 11206248}, 2018.
\bibitem{tensorsig} E. Aune. \emph{tensorsig}, version 0.1.0. Software, Zenodo, 2026. \href{https://doi.org/10.5281/zenodo.23886748}{doi:10.5281/zenodo.23886748}.
\end{thebibliography}

\end{footnotesize}
\end{document}
